\begin{acknowledgements}

行文至此，落墨为终。西长安街618号的故事也到了尾声，四季轮转间，银杏叶落了又生，晨光与晚霞交替了一千四百余次，来时是秋风送爽，走时是盛夏蝉鸣，中间隔着整个青春。总觉得毕业遥遥无期，却也转瞬到了落笔于此的时刻。四年前，我怀揣着青涩与憧憬步入校园；四年后，拨穗在即，这段充满挑战与温情的求学之路即将画上句点。值此论文完稿之际，心中百感交集，万般不舍皆化作笔尖的诚挚谢意。

一朝沐杏雨，一生念师恩。我要向我的指导老师致以最深切的谢意。从最初选题的迷茫、架构的构思，到代码的反复调试、论文的逐句修改，您的悉心指导与包容鼓励，如同引航的灯塔，帮我拨开迷雾，照亮前行之路。您严谨治学的态度、深厚的专业造诣与无私的关怀，使我受益终身，不仅教我求知，更育我成长。同时，也要感谢通信工程专业所有授课老师的辛勤培育，是你们传授的扎实理论为我这几年的探索奠定了坚实的基石。

山水一程，三生有幸。感谢朝夕相伴的同窗与挚友们。那些在自习室并肩奋斗的深夜、在实验室里为攻克技术难关而热烈讨论的午后，以及在迷茫失落时给予的温暖拥抱，都拼凑成了我青春里最灿烂的拼图。特别感谢那些在毕业设计期间，主动将自家长辈手中各品牌手机借给我做兼容性测试的朋友们，你们的无私信任与慷慨相助，让我的课题能够扎根于真实而广泛的设备之上。愿我们无论走向何方，都能跃入人海，各呈璀璨。

树高千尺不忘根深沃土。最想感谢的是我的家人，谢谢你们无条件的爱与信任，成为我人生路上最坚实的后盾。本课题以适老化软件系统为研究对象，最初的灵感便源自对长辈使用智能手机时重重困境的观察。写下这些文字时，我常会想起十八岁那年，站在讲台上向大家说出目标院校的自己——那时的我眼神清亮，大声说自己坚信“科技改变世界，科技造福人类”。如今，我终于用这四年的所学，为老龄化社会的数字包容递出了一份温热的答卷。回首来路，那个笨拙前行的少年没有食言，他守护住了最初的梦想与初心。

最后，感谢那个从不轻言放弃的自己，在无数次自我否定与重塑中，依然选择勇敢地走下去。二十二岁的天空开阔而明朗，我知道，我不能一直仰望星空，我要走入风雨，去真正拥抱和改变这个世界了。

已渡万重山，便不惧风浪。落幕的是学生时代，而非我的人生。这一站是终点，亦是起点，未来仍有无限可能。谨以此篇，献给我最热烈而真挚的青春。

\end{acknowledgements}
